\chapter{Տվյալների ֆորմատ և մոդել}
\label{ch:data}

FPGA/PS տվյալները սերվեր ուղարկելուց առաջ ձևափոխվում են կառուցվածքային ֆորմատի.
\begin{verbatim}
{ "device": "AUP_ZU3_01", "experiment_id": 12, "sensor": "vibration",
  "timestamp": 1789063200, "sequence": 1524,
  "raw_value": 728, "processed_value": 681, "status": "normal" }
\end{verbatim}
PS-ը տվյալները ուղարկում է API-ին, օրինակ՝ \texttt{POST /api/v1/data}։

Տվյալների բազան պետք է աջակցի \textbf{փորձերի (experiments)} հասկացությունը.
\begin{itemize}
  \item \textbf{Experiment} — \texttt{experiment\_id}, \texttt{device\_id}, \texttt{configuration}, \texttt{start\_time}, \texttt{end\_time}։
  \item \textbf{Sensor} — \texttt{sensor\_id}, \texttt{sensor\_type}, \texttt{sample\_rate}, \texttt{unit}։
  \item \textbf{Sample} — \texttt{experiment\_id}, \texttt{sensor\_id}, \texttt{timestamp}, \texttt{sequence\_number}, \texttt{raw\_value}, \texttt{processed\_value}, \texttt{status}։
\end{itemize}
Այս մոդելը թույլ է տալիս համեմատել փորձերը (օր.՝ Experiment~\#12 vs \#18)։
